\subsection{Khung so sánh bốn họ cổng}
\label{muc:hocong}

Ba định lý trên so hai cơ sở tính của \emph{cùng một} cơ chế. Các công trình gần nhất lại dùng
những ràng buộc lời hứa có hình thức khác: Schulz và Pfeiffer \cite{schulz2026} đặt một hạn chót
tính từ lời hứa đầu tiên; một cửa sổ quanh lời hứa đầu gần với ``hộp tĩnh'' gồm cửa sổ thời gian
và thời gian đi tối đa của Sharif Azadeh và cộng sự \cite{sharifazadeh2022}. Mục này đặt tất cả
vào cùng một khung để so được, và rút ra kết quả trung tâm của chương.

\paragraph{Đại lượng có dấu và đại lượng cộng dồn.} Xét một hành khách và tọa độ giờ trả. Tại mỗi
lần công bố $t$ đặt
\begin{equation}
x_t = e_t - v_t, \qquad c_t = p_t - e_t, \qquad z_t = p_t - v_t,
\end{equation}
khi đó $d_t = |x_t|$, $A_t = |c_t|$, $B_t = |z_t|$, và đồng nhất \emph{có dấu} $z_t = x_t + c_t$
luôn đúng --- bất đẳng thức \eqref{eq:phanra} chỉ xuất hiện khi lấy trị tuyệt đối. Gọi $p_0$ là
lời hứa \emph{đầu tiên} của hành khách. Cộng dồn trên cùng lưới công bố:
\begin{equation}
C = \sum_t A_t, \quad X = \sum_t d_t, \quad V = \sum_t B_t, \quad
M = \max_t |p_t - p_0|, \quad M^{+} = \max_t \max(0,\, p_t - p_0).
\end{equation}
$M^{+}$ là phần \emph{muộn} cực đại so với lời hứa đầu --- đúng đại lượng mà một hạn chót kiểu
Schulz--Pfeiffer chặn.

\begin{menhde}[Chuỗi đại lượng]
\label{md:chuoi}
Nếu chuỗi công bố liên tục, tức mỗi lần công bố so với lần ngay trước nó ($v_t = p_{t-1}$), thì
\[
\max(0,\, p_T - p_0) \;\le\; M^{+} \;\le\; M \;\le\; V \;\le\; C + X .
\]
\end{menhde}

\begin{proof}
Bất đẳng thức đầu và thứ hai là $\max(0,\cdot) \le |\cdot|$ lấy cực đại trên tiền tố. Với chuỗi
liên tục, $p_t - p_0 = \sum_{s \le t} z_s$, nên $|p_t - p_0| \le \sum_{s \le t} |z_s| \le V$.
Cuối cùng $|z_t| \le |c_t| + |x_t|$ ở từng lần công bố, cộng dồn lại được $V \le C + X$.
\end{proof}

Bản thân chuỗi bất đẳng thức này là kiến thức cơ bản, \emph{không} phải đóng góp; nó chỉ là công
cụ cho định lý dưới đây.

\paragraph{Bốn họ cổng.} Bảng \ref{bang:hocong} xếp bốn cách ràng buộc lời hứa theo đại lượng mà
chúng kiểm ở mỗi lần công bố.

\begin{table}[htbp]
\centering
\caption{Bốn họ cổng cam kết với trần $\kappa$}
\label{bang:hocong}
\small
\begin{tabular}{@{}llll@{}}
\toprule
\textbf{Họ} & \textbf{Kiểm ở mỗi lần công bố} & \textbf{Tên gọi} & \textbf{Ví dụ}\\
\midrule
$C$ & $\sum_t A_t \le \kappa$ & theo hành vi & cơ sở hành vi của khóa luận\\
$M^{+}$ & $p_t - p_0 \le \kappa$ & hạn chót một chiều & Schulz và Pfeiffer \cite{schulz2026}\\
$M$ & $|p_t - p_0| \le \kappa$ & hạn chót hai chiều & cửa sổ quanh lời hứa đầu\\
$V$ & $\sum_t B_t \le \kappa$ & theo kết quả & cơ sở kết quả của khóa luận\\
\bottomrule
\end{tabular}
\end{table}

\paragraph{Phạm vi của định lý và giả thiết (N).} Định lý dưới đây chỉ nói về \emph{phần cam kết}
của bộ lọc. Tập phương án của một xe vẫn có thể rỗng, hoặc lần chạy vẫn có thể dừng, vì những ràng
buộc \emph{không} thuộc cơ chế cam kết: cấu trúc lộ trình, sức chứa, và nhất là bất biến vòng đời
khi xe đón muộn thật. Điều này đã xảy ra: dưới trôi dạt đo thật của cao điểm sáng (mục
\ref{muc:troidat}), chính cơ sở hành vi với trần 30 giây dừng ở 11 trên 20 ô, và cả 11 lần dừng
đều do bất biến cửa sổ đón của tầng miền (mục \ref{muc:batbien}) chứ không do cam kết. Cơ chế cụ
thể: bộ sinh phương án miễn kiểm cửa sổ dịch vụ \emph{chỉ} cho phương án không làm gì
(\texttt{ServiceQualityAllowance}; mọi phương án có thay đổi đều bị kiểm chặt), nên phương án này
không bị loại vì cửa sổ; lần chạy dừng muộn hơn, khi xe thật sự tới điểm đón sau giờ đón muộn nhất
và thao tác lên xe bị tầng miền từ chối (\texttt{RideRequest.Board}). Gọi \textbf{(N)} là giả
thiết: \emph{tại kỷ nguyên đang xét, phương án không làm gì của xe đã qua mọi kiểm tra không thuộc
cam kết.}

\paragraph{Quỹ đạo dự phòng.} Xét quỹ đạo mà ở mọi kỷ nguyên xe giữ nguyên tuyến. Trên đó
$p_t = e_t$, nên $c_t = 0$ với mọi $t$, và
\[
C_{\text{noop}} = 0, \qquad z_t = x_t, \qquad V_{\text{noop}} = X .
\]

\begin{dinhly}[Thứ tự khả thi của các họ cổng]
\label{dl:hocong}
Giả sử (N). Xét một hành khách.
\begin{enumerate}[label=(\alph*),leftmargin=1cm,itemsep=2pt,topsep=3pt]
\item \emph{Mọi quỹ đạo.} Với mọi $\kappa \ge 0$ và trên mọi quỹ đạo, cổng $C$ không bao giờ loại
phương án không làm gì. Do đó cổng hành vi \emph{không bao giờ là nguyên nhân} làm rỗng tập phương
án của một xe.
\item \emph{Quỹ đạo dự phòng.} Trên quỹ đạo dự phòng với chuỗi công bố liên tục, ngưỡng nhỏ nhất
bảo đảm phương án không làm gì không bị loại ở mọi kỷ nguyên là
\[
\kappa^{*}_{C} = 0 \;\le\; \kappa^{*}_{M^{+}} = M^{+}_{\text{noop}} \;\le\;
\kappa^{*}_{M} = M_{\text{noop}} \;\le\; \kappa^{*}_{V} = X ,
\]
tức ba họ còn lại có thể loại phương án không làm gì, và ngưỡng an toàn của chúng tăng theo độ
trôi dạt.
\end{enumerate}
\end{dinhly}

\begin{proof}
(a) Sổ cái trước kỷ nguyên $t$ thỏa $C_{t-1} \le \kappa$ vì lần công bố trước đã qua cổng. Phương
án không làm gì có $c_t = 0$ nên $C_t = C_{t-1} \le \kappa$: cổng không loại nó, và theo (N) nó
cũng qua mọi kiểm tra khác, nên tập phương án không bị làm rỗng \emph{vì cổng}. (b) Phương án không
làm gì bị loại ngay tại lần công bố đầu tiên mà đại lượng được kiểm vượt $\kappa$, nên ngưỡng nhỏ
nhất đúng bằng giá trị lớn nhất của đại lượng đó dọc quỹ đạo dự phòng (với $V$ là giá trị cuối, vì
$V$ cộng dồn đơn điệu). Thứ tự giữa các ngưỡng là Mệnh đề \ref{md:chuoi} áp cho quỹ đạo có
$C = 0$; bước $M \le V$ cần chuỗi công bố liên tục.
\end{proof}

Phần (a) chính là Định lý \ref{dl:hanhvi} phát biểu lại trong khung chung; phần mới là việc đặt ba họ còn lại
--- trong đó có cơ chế của văn liệu --- vào cùng thước so. Chứng minh là \emph{sơ cấp}, và khóa
luận nói thẳng như vậy: giá trị của định lý không nằm ở độ khó mà ở hai chỗ --- đặt các họ cổng
trong văn liệu vào một khung để so được, và kiểm chứng rằng ranh giới nó vạch ra đúng với hệ thật.

Câu phát biểu chính xác là: \emph{cổng hành vi không bao giờ là nguyên nhân khiến tập phương án
rỗng}. Câu ``hệ thống dùng cổng hành vi không bao giờ kẹt'' là \emph{sai} --- 11 trên 20 ô ở trên
là phản ví dụ.

\begin{hequa}[Khi thời gian đi tất định]
\label{hq:tatdinh}
Nếu thời điểm tới của tuyến giữ nguyên không trôi --- thời gian di chuyển và thời gian thực thi
đều tất định --- thì $x_t \equiv 0$, do đó $M^{+}_{\text{noop}} = M_{\text{noop}} = X = 0$ và cả
bốn họ đều không bao giờ loại phương án không làm gì, với mọi $\kappa \ge 0$.
\end{hequa}

Trong mô hình của Schulz và Pfeiffer, thời gian đi giữa hai điểm dừng là một giá trị cố định cho
mỗi cặp (mục 3.1 của bài), và bài khẳng định khách đã được nhận thì không bao giờ bị từ chối về
sau. Hệ quả \ref{hq:tatdinh} nhất quán với điều đó: trong một mô hình thời gian đi tất định, hạn
chót tính từ lời hứa đầu không thể tự nó làm rỗng tập phương án. Tính vô nghiệm chỉ xuất hiện khi
lời hứa của tuyến giữ nguyên bị trôi --- điều nằm ngoài mô hình của họ.

\begin{hequa}[Phụ thuộc chiều trôi]
\label{hq:chieutroi}
Trên quỹ đạo dự phòng với chuỗi công bố liên tục, $p_t - p_0 = \sum_{s \le t} x_s$, nên
\[
M^{+}_{\text{noop}} = \max_t \Big(\sum_{s \le t} x_s\Big)^{+}, \qquad
M_{\text{noop}} = \max_t \Big|\sum_{s \le t} x_s\Big|, \qquad
V_{\text{noop}} = \sum_t |x_t| .
\]
Nếu trôi \emph{bất lợi} (mọi $x_s \ge 0$, tắc đường tăng) thì
$M^{+}_{\text{noop}} = M_{\text{noop}} = V_{\text{noop}}$: ba họ $M^{+}$, $M$, $V$ cùng vô nghiệm ở
cùng một ngưỡng. Nếu trôi \emph{thuận lợi} (mọi $x_s \le 0$, tắc đường giảm) thì
$M^{+}_{\text{noop}} = 0$: hạn chót một chiều an toàn với mọi $\kappa$, trong khi
$M_{\text{noop}} = V_{\text{noop}}$ vẫn có thể lớn, nên hạn chót hai chiều và cổng kết quả vẫn có
thể vô nghiệm. Cổng hành vi an toàn ở cả hai chiều.
\end{hequa}

Hệ quả này là chỗ làm cho định lý \emph{kiểm được} chứ không chỉ đúng. Một nhận định thô kiểu
``hạn chót hỏng khi có trôi dạt'' không phân biệt được hai chiều; Định lý \ref{dl:hocong} cùng Hệ
quả \ref{hq:chieutroi} đoán \emph{họ nào} hỏng theo \emph{chiều} trôi. Trên hệ thật $x_s$ không
thuần một dấu --- nó còn gồm tiến độ thực và độ trễ thực thi của xe --- nên các đẳng thức trên chỉ
đúng gần đúng, và mọi dự đoán phải tính từ quỹ đạo đo được của từng khách chứ không từ công thức.

Hệ quả cũng có một mặt trái phải được báo cáo. Dưới hạn chót một chiều và trôi thuận lợi, một
phương án chèn khách làm lời hứa muộn đi ($c_t > 0$) vẫn qua cổng khi $e_t - p_0$ đủ âm --- tức nhà
điều hành được \emph{tiêu phần trôi thuận lợi} để dời khách, trong khi cổng hành vi tính đủ $c_t$.
Vì vậy trong những khung giờ tắc đường giảm, hạn chót một chiều có thể phục vụ nhiều khách hơn trần
hành vi cùng mức, đổi lại khách thấy lời hứa đi sớm rồi muộn lại. Đây là một đánh đổi thật giữa hai
họ, không phải một ưu thế một chiều.

\subsection{Khi nào chặn theo chuỗi khác chặn hai đầu}
\label{muc:dondieu}

\begin{menhde}[Bộ sinh chèn thuần]
\label{md:dondieu}
Xét một chuỗi công bố liên tục. Nếu mọi phương án được chọn đều có $p \ge e$ với mọi hành khách
cũ, tức $c_t \ge 0$ ở mọi lần công bố, thì $C = \sum_t c_t$. Nếu thêm $x_t \equiv 0$ thì
$C = V = p_T - p_0 = M^{+} = M$: trần theo chuỗi, trần theo hai đầu và trần theo hành vi là một.
\end{menhde}

\begin{proof}
$c_t \ge 0$ cho $A_t = c_t$. Khi $x_t \equiv 0$ thì $z_t = c_t \ge 0$; với chuỗi liên tục
$p_t - p_0 = \sum_{s \le t} z_s$, nên $p_t$ đơn điệu không giảm, $p_T - p_0 = \sum_t z_t = V = C$,
và cực đại của $p_t - p_0$ đạt ở lần công bố cuối.
\end{proof}

Giả thiết của mệnh đề đúng trên hệ này ở mức đo được: dấu của $p - e$ dương ở cả $479$ lần công bố
có $A > 0$ của B1 (mục \ref{muc:phanra}), vì bộ sinh chỉ chèn thêm. Trên panel tĩnh, số khách vượt
một ngưỡng tính theo hai đầu và tính theo thành phần hành vi khác nhau không quá một khách trên
$1.735$ ở mọi ngưỡng 30, 60, 120, 300 giây. Hệ quả thực hành: phát biểu ``chặn tổng biến thiên mạnh
hơn chặn hai đầu'' đúng về toán nhưng \emph{rỗng} trên hệ này; nó chỉ có hiệu lực khi bộ sinh có thể
kéo lời hứa sớm lên, tức khi có sắp xếp lại. Khóa luận vì vậy \emph{không} dùng tổng biến thiên làm
điểm phân biệt chính.

\subsection{Các phản ví dụ}

Bảng \ref{bang:phanvidu} gom các tình huống nhỏ minh họa ranh giới của từng đại lượng. Đơn vị là
giây; năm dòng đầu là ví dụ dựng tay, dòng cuối là một lần công bố thật.

\begin{table}[htbp]
\centering
\caption{Phản ví dụ cho các đại lượng của khung bốn họ cổng}
\label{bang:phanvidu}
\small
\begin{tabular}{@{}llp{6.4cm}@{}}
\toprule
\textbf{Tình huống} & \textbf{Kết quả} & \textbf{Minh họa điều gì}\\
\midrule
$v = 600,\ e = 600,\ p = 660$ & $x = 0,\ c = 60,\ z = 60$ & Quyết định thuần gây dời.\\
$v = 600,\ e = 660,\ p = 660$ & $x = 60,\ c = 0,\ z = 60$ & Giữ tuyến mà lời hứa vẫn xấu đi: trần
hành vi \emph{không} chặn cái khách thấy.\\
$v = 600,\ e = 720,\ p = 600$ & $C = 120,\ V = 0$ & Sửa để giữ lời hứa bị trần hành vi phạt; chưa
từng xảy ra trên panel ($0/479$).\\
$p_0 = 600$; công bố $660$, $600$ & $E = 0,\ V = 120$ & Thước hai đầu mù với dao động.\\
$p_0 = 600$; công bố $660$, $720$ & $E = 120,\ V = 120$ & Đơn điệu thì chuỗi bằng hai đầu.\\
Ca thật (dưới) & $A = 90{,}244;\ d = 0{,}690;\ B = 89{,}554$ & Trôi thuận chiều hấp thụ một phần
cú đẩy; hụt đúng $2d = 1{,}380$.\\
\bottomrule
\end{tabular}
\end{table}

Dòng cuối là một lần công bố thật ở ô 14-r1 của nhánh B1, lúc thời điểm mô phỏng 00:06:36:
$v = $ 00:15:09,943, $e = $ 00:15:09,253 và $p = $ 00:16:39,497 (giờ:phút:giây mô phỏng). Tuyến giữ
nguyên đã sớm hơn lời hứa $0{,}690$ giây, rồi quyết định đẩy giờ trả muộn thêm $90{,}244$ giây;
khách chỉ thấy $89{,}554$ giây. Đây cũng là một lần (H) hỏng ($A > B$), đúng loại ghi nhận ở bảng
\ref{bang:dieukien-h}.

Các kiểm chứng thực nghiệm của khung này nằm ở mục \ref{muc:troidat}, \ref{muc:haunghiem} và
\ref{muc:hoito}. Ở thời điểm viết, chúng là \emph{chứng minh kèm giải thích hậu nghiệm} trên dữ
liệu có sẵn, chưa phải thí nghiệm xác nhận có tiền đăng ký.
